% Zusammenfassung
\section{Zusammenfassung der wichtigsten Ergebnisse}

Zur Beantwortung der Problemstellung, inwiefern Dashbaords geeignet und sinnvoll sind, als Modernisierungsmethode eingesetzt zu werden, wurde hierfür eine Nutzerstudie bestehend aus einem Usability-Test und einer dazugehörigen Nutzerbefragung durchgeführt. Während die Nutzerbefragung auf den Prinzipien der ISO 9241-110 basiert, wurden innerhalb des Usability-Tests die nachträglich abgefragten Eigenschaften eingebaut. Auf dieser Grundlage konnten die beide Systeme anhand der drei Usability-Dimensionen Effektivität, Effizienz und Nutzerzufriedenheit verglichen werden.

%Anhand dieser Forschungen konnten einige Möglichkeiten festgestellt werden.

Mithilfe der gewählten Dash\-board"=Lösung ist ein zeitlich effizienteres Monitoring möglich als zuvor. Zwar wies der Usability-Test an dem Dash\-board eine höhere Fehlerquote auf als an dem Altsystem, jedoch liegt dies an der exzellenten Kenntnis der Nutzer des Altsystems sowie der Tatsache, dass die neue Lösung erst mit der Teilnahme an der Studie vorgestellt wurde. Allerdings konnte das Dash\-board in den entscheidenden Testbereichen durch seine verbesserte Effektivität überzeugen und somit eine um fast 50\% bessere Effizienz vorweisen.

Anhand der regen Teilnahme an der Feedbackrunde war festzustellen, dass durchaus Verbesserungspotenzial an dem neuen Prototypen besteht. Jedoch wird damit auch Interesse auf Nutzerbasis gezeigt, dieses Projekt weiter zu fördern. Insgesamt sind die Nutzer auch viel zufriedener mit der Umsetzung und Nutzbarkeit des Dash\-boards als mit der gewachsenen Altlösung.

% TODO: werden folgende Fragen beantwortet? - wenn nein, dann noch machen
% Wie sinnvoll ist die Einführung eines modernen Dashboards im Hinblick auf das Ziel, den Legacy-Testpoolmonitoring-Prozess effizienter und nutzerfreundlicher zu gestalten?
%Zur Beantwortung sollen anhand der nachfolgenden Methodik (s. Abschnitt 3.4) folgende Teilaspekte geklärt werden:
%• Wie weit kann die Usability mithilfe des Dashboards verbessert werden?
%• Wie zufrieden sind die Nutzer mit dem Dashboard? Ist ein Umstieg auf das neue System denkbar?
%• Wie unterscheidet sich das Dashboard anhand der Intuitivität vom alten System? Ist es auch für neue Nutzer ohne Vorwissen geeignet?
%• Wie wirksam ist die Nutzung des Dashboards in Bezug auf die Effizienz des Monitoring-Prozesses?
%• Beeinflusst die Nutzung des Dashboards die Effektivität des Monitoring-Prozesses?
